% The Tate module and the Galois representation, stacked vertically.  The torsion
% E[l^n]=(Z/l^n)^2 forms an inverse system under x l; its limit is T_l E = Z_l^2, on which
% the absolute Galois group acts, giving  rho : Gal(Qbar/Q) -> GL_2(Z_l).
\begin{tikzpicture}[>=Latex,
    tb/.style={rectangle, rounded corners=5pt, draw=formteal!70!black, line width=1pt,
      fill=formteal!8, inner sep=7pt, align=center}]
  % inverse system (top row)
  \node[tb] (e1) at (0,0)   {$E[\ell]$};
  \node[tb] (e2) at (2.6,0) {$E[\ell^2]$};
  \node[tb] (e3) at (5.2,0) {$E[\ell^3]$};
  \node[font=\Large] at (6.9,0) {$\cdots$};
  \draw[<-, line width=1.4pt, draw=formslate] (e1)--node[above, font=\small] {$\ell$} (e2);
  \draw[<-, line width=1.4pt, draw=formslate] (e2)--node[above, font=\small] {$\ell$} (e3);
  \draw[<-, line width=1.4pt, draw=formslate] (e3)--node[above, font=\small] {$\ell$} (6.55,0);
  \node[font=\small, formteal!55!black, anchor=north] at (2.6,-0.9)
    {each $E[\ell^n]\cong(\Z/\ell^n)^2$};
  % Tate module (middle), centred
  \node[tb, draw=formgold!84!black, fill=formgold!15] (T) at (3.0,-3.9)
    {\textbf{Tate module}\\[2pt]$T_\ell E=\varprojlim_n E[\ell^n]\cong\Z_\ell^{2}$};
  \draw[->, dotted, line width=1pt, draw=formslate] (e2.south) to[bend right=10] (T.north);
  % Galois representation (bottom), centred
  \node[tb, draw=formrust!76!black, fill=formrust!10] (G) at (3.0,-7.4)
    {\textbf{Galois representation}\\[2pt]$\rho:\operatorname{Gal}(\overline{\Q}/\Q)\to\mathrm{GL}_2(\Z_\ell)$};
  \draw[->, line width=1.5pt, draw=formslate] (T)--node[right=1pt, font=\small] {Galois acts} (G);
\end{tikzpicture}
